\subsection{Comparaison de deux corps purs}
\Q Le méthanol gazeux est assimilé à un gaz parfait. Calculer la variation du potentiel chimique de \ce{H3C-OH(g)} qui, à la température de \SI{298}{\kelvin} et à partir de la pression $P=P^\circ= 1 ~ \si{\bar}$, subit une augmentation
de pression $\Delta P = 1 ~ \si{\milli\bar}$.

\Q L'activité d’un gaz parfait est multipliée par $10$, à température constante. Calculer, à \SI{298}{\kelvin}, la variation correspondante de son potentiel chimique.

\Q Calculer, à 298 K, les pressions de vapeur saturantes du tétrachlorométhane CC liquide et de l’eau liquide.

\begin{center}
\begin{tabular}{ccc}
\toprule
Constituants / Etats de la matière & liquide & gaz \\
\midrule
\ce{H2O} & - 237.2 & -228.6 \\
\ce{CCl4} & -68.74 & -64.22 \\
\bottomrule
\end{tabular}
\captionof{table}{Potentiels chimiques standard des corps purs, à \SI{298}{\kelvin}, dans différents états de la matière}
\end{center}